\subsection{An exact Fourier criterion on an odd quotient}
\label{sec:odd-quotient-fourier}

\paragraph{Finding A49 (all conductors of a prescribed odd quotient).}
Let $C/k$ be a smooth projective geometrically connected genus-two
curve over $k=\mathbf F_q$ of characteristic $p$, and let $P\in C(k)$.
Put $G=J_C(k)$, $N=\#C(k)$ and $\kappa_P=[K_C-2P]$.
Let
\[
 b:G\twoheadrightarrow H,\qquad
 n=\exp H>1,\qquad n\text{ odd},\qquad p\nmid n
 \tag{OQ.1}
\]
be a finite abelian quotient satisfying
\[
 b(\kappa_P)=0,\qquad
       \langle p,-1\rangle=(\mathbf Z/n)^\times.
 \tag{OQ.2}
\]
For $h\in H$, let $[h]$ denote the corresponding basis element of
$\mathbf F_p[H]$, and define
\[
 F_H=[0]+\sum_{R\in C(k)}[b([R-P])],\qquad
 U_H=\frac{N+1}{|H|}\sum_{h\in H}[h].
 \tag{OQ.3}
\]
The fraction is taken in $\mathbf F_p$; its denominator is invertible.
Then the following are equivalent:
\begin{enumerate}
\item $F_H\ne U_H$ in $\mathbf F_p[H]$.
\item There is a cyclic quotient $c:H\twoheadrightarrow\mathbf Z/m$,
      with $m>1$ dividing $n$, whose constant $P$-split Lang cover
      $D\to C$ over $k$ has actual primitive conductor-$m$ factor
      $\mathcal P_m$ satisfying
      \[
                  p\nmid\#\mathcal P_m(k).
       \tag{OQ.4}
      \]
\end{enumerate}
The cover in the second assertion is geometrically connected.
When these equivalent conditions hold, A46 separates the point section of the prescribed $P$
from that of every distinct point of $C(k)$.
The conclusion concerns the primitive factor of this cover;
proper-conductor factors may have rational $p$-primary torsion.

In particular, the sufficient numerical conditions
\[
                 p\nmid(N+1),\qquad |H|>N
 \tag{OQ.5}
\]
imply (OQ.4) for some such cyclic quotient. In characteristic two,
the first condition is simply that $N$ is even. It therefore holds
for every genus-two curve of two-rank one by A26.
At a rational Weierstrass point $P$, the canonical condition in
(OQ.2) holds for every quotient, since $K_C\sim2P$.

\paragraph{Recovered inputs and the scope of the refinement.}
The Fourier identity and empty-coset argument are already present
in recovered successor A.10 and A.35, and A26 isolates their
even-$N$ consequence for a prime quotient on which $p$ is primitive
\cite[A.10, A.35]{ASrun}. The genus-two character polynomial is
recovered A1 B.2; its canonical-class symmetry is the one used in
recovered successor A.26 and A29 \cite[B.2]{A1run}\cite[A.26]{ASrun}.
Here the quotient may have higher prime powers or several prime
divisors. Condition (OQ.2) uses the canonical symmetry to join the
$p$-power and inverse character orbits, and A46 permits selecting
only the resulting primitive conductor. The proof below establishes
the exact equivalence and its original-field pointing; it makes no
priority claim for the underlying Fourier or Lang constructions.

\begin{proof}
\emph{1. Every cyclic quotient gives a cover over the original field.}
Fix a cyclic quotient $c:H\to A=\mathbf Z/m$, with $m>1$, and put
$a=c\circ b$ and $K=\ker a\subset G$.
Let $\Phi$ be $q$-power Frobenius on $J=J_C$. The Lang isogeny
$\Phi-1$ is etale with constant kernel $G$, and factors over $k$ as
\[
             J\longrightarrow J/K\xrightarrow{\psi_a}J,
                    \qquad \ker\psi_a=G/K=A.
 \tag{OQ.6}
\]
Pull it back by $i_P:C\to J$, $R\mapsto[R-P]$.
The result $D_a=C\times_J(J/K)$ is an etale cover with constant
deck group $A$. The Abel--Jacobi connected-pullback theorem
\cite[Proposition~9.1]{MilneJV}, as used in A12, makes it
geometrically connected. Its fibre over $P$ is the constant kernel
of $\psi_a$, containing the origin. Thus it is $P$-split over $k$
itself. No extension to rationalize a point or deck transformation
has been made.

\emph{2. The exact primitive rational test.}
For a nontrivial characteristic-zero character $\chi$ of $A$, use
the Lang Frobenius convention of recovered A1 B.2. Its exact
unramified genus-two polynomial and middle-coefficient symmetry are
\[
\begin{aligned}
 L_P(\chi,T)&=1+S_\chi T+q\chi(a(\kappa_P))T^2,\\
 S_\chi&=\sum_{R\in C(k)}\chi(a([R-P])),\\
 S_\chi&=\chi(a(\kappa_P))S_{\chi^{-1}}.
\end{aligned}
 \tag{OQ.7}
\]
These are precisely the Riemann--Roch identities proved there
and in (PC.6)--(PC.7). Reversing the Lang convention inverts all
characters and changes none of the tests below. Here
$a(\kappa_P)=0$ by (OQ.2), so $S_\chi=S_{\chi^{-1}}$.

Define $\mathcal P_m$ to be the actual image of the integral
canonical projector $m e_m$ on $J_{D_a}$, as in A46; $e_m$ is
the rational idempotent supported on the characters of exact
order $m$. It is defined over $k$. The deck-character decomposition
of the zeta numerator, or equivalently its action on a Tate module
of auxiliary characteristic prime to $pm$, gives the exact identity
\[
       \#\mathcal P_m(k)
       =\prod_{\substack{\chi\in\widehat A\\\operatorname{ord}\chi=m}}
                         L_P(\chi,1).
 \tag{OQ.8}
\]
The image has exactly these rational character summands; finite-field
point counts depend on the Frobenius polynomial, so this equality
uses the actual image and does not require a chosen isogenous
representative. All factors on the right are algebraic integers.

Fix a prime above $p$ in a cyclotomic field containing the $n$th
roots of unity. Reduction identifies their character group with
the $n$-torsion roots of unity in $\overline{\mathbf F}_p$, since
$p\nmid n$. By (OQ.7),
\[
        \overline{L_P(\chi,1)}
          =1+\sum_{R\in C(k)}\bar\chi(a([R-P]))
          =(\bar\chi\circ c)(F_H).
 \tag{OQ.9}
\]
Consequently (OQ.4) is equivalent to nonvanishing of (OQ.9)
for every faithful character of $A$. Indeed the product (OQ.8)
is a rational integer, and its reduction is nonzero exactly
when every factor has nonzero reduction at this fixed prime.

\emph{3. One character controls its entire conductor.}
Reduction of units $(\mathbf Z/n)^\times\to(\mathbf Z/m)^\times$
is surjective for every $m\mid n$. Thus (OQ.2) implies
\[
             \langle p,-1\rangle=(\mathbf Z/m)^\times.
 \tag{OQ.10}
\]
For any character $\eta:H\to\overline{\mathbf F}_p^\times$,
the coefficients of $F_H$ lie in $\mathbf F_p$, giving
\[
                 \eta^p(F_H)=\eta(F_H)^p.
 \tag{OQ.11}
\]
The hyperelliptic involution is defined over $k$ and satisfies
$[\iota R-P]=\kappa_P-[R-P]$.
As $b(\kappa_P)=0$, its permutation of $C(k)$ also gives
\[
                         \eta^{-1}(F_H)=\eta(F_H).
 \tag{OQ.12}
\]
Hence vanishing, and also nonvanishing, is preserved by the
$p$-power and inverse character operations. If $\eta$ has
exact order $m$, all generators of its cyclic character subgroup
are reached using (OQ.10). Thus a single nonzero value supplies
nonzero values for every faithful character of the quotient
$H/\ker\eta$ of order $m$.

\emph{4. Fourier inversion and the numerical corollary.}
The finite Fourier transform on $H$ is invertible over
$\overline{\mathbf F}_p$, because $p\nmid|H|$.
The trivial character evaluates $F_H$ to $N+1$.
All its nontrivial character values vanish if and only if
\[
                   F_H=\frac{N+1}{|H|}\sum_{h\in H}[h]=U_H.
 \tag{OQ.13}
\]
If $F_H\ne U_H$, choose a nontrivial character $\eta$ with
$\eta(F_H)\ne0$. Its image has order $m>1$ dividing $n$.
Step 3 and (OQ.9) make every faithful factor of its cyclic
quotient nonzero; Step 2 proves (OQ.4), and Step 1 supplies
the required original-field cover. Conversely, a quotient
satisfying (OQ.4) has a faithful character with nonzero value
on $F_H$, so (OQ.13) cannot hold. This proves the equivalence.

Under (OQ.5), the coefficient $(N+1)/|H|$ is nonzero.
There are at most $N$ occupied Abel-image cosets, and the zero
coset is occupied by $P$. Since $|H|>N$, an empty nonzero coset
has coefficient zero in $F_H$ and a nonzero coefficient in
$U_H$. Thus $F_H\ne U_H$. Finally (OQ.4) says
$\mathcal P_m(k)[p^\infty]=0$, so the point-separation
conclusion follows directly from A46, with its odd cyclic
degree and prescribed rational lift hypotheses all satisfied.
\end{proof}

\paragraph{Two retained comparisons over $\mathbf F_4$.}
On A27's curve
$y^2+y=ux^3+x+u/x$, $u^2+u+1=0$, take either rational
Weierstrass point $W$. The retained counts are
$\#C(\mathbf F_4)=6$ and $\#J_C(\mathbf F_4)=28$.
Its seven-primary quotient is $H=\mathbf Z/7$.
The powers of two modulo seven are $\{1,2,4\}$, and their
negatives give the other three units. Thus (OQ.2) and (OQ.5)
hold and A49 gives a constant $W$-split degree-seven primitive
rational witness over $\mathbf F_4$. The primitive-root
hypothesis of A26 fails at seven. A31 already supplies a
geometric degree-seven witness for this same curve; the
comparison concerns the reach of the sufficient criteria.

For the quadratic twist used in A45, namely
\[
          C:\ Y^2+X^2Y=uX^5+uX^4+X^3+uX,
                     \qquad W=\infty,
 \tag{OQ.14}
\]
the retained data (ND.12)--(ND.19) give $N=4$,
$\#J_C(\mathbf F_4)=18$, and a rational class of exact
order nine. Hence its full three-primary quotient is
$H=\mathbf Z/9$. The powers of two modulo nine are
$1,2,4,8,7,5$, so the unit condition holds; $\kappa_W=0$
and $9>4$. The only possible nontrivial cyclic conductors
are three and nine. The unique degree-three quotient,
up to an automorphism of its deck group, has
$L_\chi(T)=1+T+4T^2$ for each nontrivial character.
Its primitive factor has order $6^2$, and therefore fails
(OQ.4). A49 consequently supplies a degree-nine quotient
whose primitive conductor-nine factor has odd rational order.
Here the prime-level quotient has size $3\le4$, while the
full primary quotient has size $9>4$.

One can check that last odd order exactly from the same retained
point data, without a new curve search. The two Weierstrass
points have label zero in $H$; the other two points have labels
$a,-a$ with $3\nmid a$, because their reductions modulo three
are $1,2$ in (ND.17)--(ND.19). For a primitive ninth root $z$,
the faithful values at $T=1$ are $7+z^j+z^{-j}$, for
$j\in(\mathbf Z/9)^\times$. The three distinct values
$s=z^j+z^{-j}$ satisfy $s^3-3s+1=0$, each occurring twice.
Therefore
\[
 \#\mathcal P_9(\mathbf F_4)
     =\left(7^3-3\cdot7-1\right)^2
     =321^2=103041.
 \tag{OQ.15}
\]
The proper conductor-three factor still has even rational order;
no assertion about the full relative class number being odd
is made. The accompanying finite algebra checker verifies the
unit orbits and the ninth-root product. Its input curve counts
and labels are the separately retained A27/A45 certificates,
not a new enumeration in that checker.

\paragraph{The remaining supply question.}
This is a test on a quotient which must actually exist over the
displayed field. A49 does not guarantee a quotient satisfying
(OQ.2), nonuniformity in (OQ.3), or the sufficient bound (OQ.5)
for every curve and prescribed point. In particular a nonzero
Fourier value without (OQ.2) need not control a whole faithful
conductor. The canonical-kernel requirement is also explicit;
A29 addresses a different situation, where a suitable primary
canonical component is nonzero. The two comparisons above
already have other witnesses in this run, so they do not prove
that the full previously unresolved universal witness gate has
shrunk. No consequence about this geometric supply is inferred
from Artin-prime density or from number-field GRH.
